\subsection*{Recurrent selection of source spans and writing actions}
The preceding visual reader encodes a joint pictured question and passage. We retained its full contextual encoder, visual regression and glyph metric. For each passage cell we cached its contextual state, inherited endpoint scores, relative source position, and the presence and last relative position of an identical bitmap in the question. These features carry no Unicode or historical stage identifiers. A linear projection and a seven-cell convolution give a contextual vector per source cell. The final classification-marker state initializes a 128-dimensional recurrent decoder.

At decision $t$ the model predicts an action: copy a contiguous source interval, emit a training-memory bitmap, or stop. The start position is scored by the inner product of the query with source keys plus the inherited start score. A selected start conditions the endpoint distribution over the following 64 positions, including the start itself. The two endpoint vectors of a copied interval, or the embedding of a written image, supply feedback for the next decision. At training time we align each reference image stream with the longest matching source interval at its current position, breaking ties in source order. Cells with no match enter a writing memory defined using the training set. Action, location, extent and writing losses are conditional cross-entropies. Greedy inference starts with empty output and stops at the learned stop decision or an explicit 200-decision limit. Only a learned stop receives exact-answer credit.

We sampled disjoint Hanziyuan families from the established source-plate corpus and retained the original raster for every historical exemplar. The training plate uses prose rather than fixed rows, varies the order of entries and stages, and sometimes requests the later form first. We used 768 training cards and 24 held-out cards. In CMRC, both orders of each selected question pair were retained, with 1,536 training pairs and 128 development pairs fitting the shared 512-cell input window. All 2,565 prepared training and 571 development DuReader records entered their respective splits. We preserved the same 64 pictured role questions as a separate regression check. We also repeated the previously available \textit{yan} (U+8A00) and \textit{zheng} (U+5F81) demonstration in both entry and request orders. This demonstration is from the prior training-family partition. The decoder receives no task label when it reads these pages.

The writer contains 1,059,840 trainable parameters and 1,331 output bitmaps learned from training. Feature preparation for both splits took 32.3 seconds with peak allocated memory of 0.79 GiB. Twelve passes took 59.4 seconds and 0.26 GiB peak allocated memory for the decoder fit with cached features. The subsequent comparison changed 6,803 pictured prompts and cached the new contexts. Each arm received the same batch memberships and 1,512 updates. Replay took 30.6 seconds and the varied fit 27.3 seconds. All times are incremental on one RTX 4090 D and exclude training the inherited contextual reader. Both independently worded tests required 48/96 complete reversals to qualify historical transfer.
